\documentclass[10pt,a4paper]{amsart}
\usepackage[margin=19mm]{geometry}
\usepackage{amsmath,amssymb,mathtools}
\usepackage[scr=rsfs,cal=euler]{mathalfa}
\usepackage{xfrac}
\usepackage[unicode,hidelinks,bookmarks=false]{hyperref}
\newcommand{\ie}{i.e.\ }
\newcommand{\cf}{cf.\ }
\newcommand{\resp}{respectively\ }
\newcommand{\Subsection}{\subsection}
\providecommand{\nobreakdash}{\nobreak\hbox{-}}
\setlength{\emergencystretch}{2em}
\multlinegap0pt
\newcommand{\F}{{\mathbb F}}
\newcommand{\Mu}{\boldsymbol \mu}
\theoremstyle{definition}
\newtheorem*{Ques}{Question}
\theoremstyle{plain}
\newtheorem{theorem}{Theorem}
\newtheorem{lemma}{Lemma}
\newtheorem{cor}{Corollary}
\newtheorem{prop}{Proposition}
\theoremstyle{definition}
\newtheorem{Def}{Definition}
\newtheorem{Ex}{Example}
\newtheorem{Rem}{Remark}
\newtheorem{Not}{Notation}
\title[Distinguishing Gauss sums]{Distinguishing Gauss sums}
\author{Moshe Adrian, Jack Diamond, Kenneth Kramer, Geo Kam-Fai Tam}
\date{Source: arXiv:2609.28211v1 (CC BY 4.0)}
\begin{document}
\maketitle

\noindent\emph{Source and licence.} The delivered work is the article shown in the title above, version arXiv:2609.28211v1 (CC BY 4.0), distributed under the Creative Commons Attribution 4.0 International licence, as recorded in the arXiv licence metadata for this exact version and shown on the abstract page saved with this item. The full licence notice, which carries the licence URL and the address of that abstract page, is the bundled file \texttt{licenses/arxiv-2609.28211-CC-BY-4.0.txt}.
\par\smallskip

\noindent\textit{Editorial scope.} Counted unit: the article's Proposition (source label Twist) (source0.tex lines 417--419, source label \texttt{Twist}): ``Assume that is prime and let and be characters on . If for all characters on , then and are Frobenius-conjugate.''. It is delivered together with the article's complete original proof (source0.tex lines 421--433, one proof environment adjacent to the statement, 13 lines), copied line for line from the article's own text. Required context retained uncounted: ConjChi (lines 247--250); EqualG (lines 256--263); f=1 (lines 277--283); gcd1 (lines 355--357). Editorial formatting only: an AMS \texttt{amsart} preamble that restates the article's own macro definitions and its own theorem declarations, and the theorem environments the retained lines use; the article's own class file is neither shipped nor input; the retained environments are renumbered sequentially in order of appearance, whereas the article prints them with its own numbering; the article's equation numbering is not reproduced and every cross-reference inside the retained text resolves to the wrapper's numbering. No citation occurs inside any retained line, so the wrapper carries no bibliography; All retained bodies are exact copies of the bundled \texttt{sources/211563/source0.tex}, so no omission receipt applies. No asset-inclusion command occurs inside any retained span and the package ships no image asset.


\section*{Retained context: ConjChi (source0.tex lines 247--250)}

\begin{equation}  \label{ConjChi}
G(\chi) G(\chi^{-1}) = \chi(-1) q
\, \text{ and } \,  \vert G(\chi) \vert_v = \sqrt{q} \,  \text{ for all } v \vert \infty.
\end{equation}


\section*{Retained context: EqualG (source0.tex lines 256--263)}

\begin{prop}  \label{EqualG}
Let $\chi = \omega^{-c}$ and $\chi' = \omega^{-c'}$.  If $p = 2$, assume $f \ge 3$.  The equality of Gauss sums $G(\chi) = G(\chi')$ holds if and only if both of the following conditions are true:  \vspace{4 pt}  

{\rm i)} \, $s_q(uc) = s_q(uc')$ for all $u \in U_{q-1}$ \vspace{5 pt} and \hspace{10 pt} 

{\rm ii)} \, $\begin{cases} \nu_q(c)  \equiv \nu_q(c') \pmod{p} &\text{if } p \text{ is odd}, \\[2 pt]
           \beta_q(c)  \equiv \beta_q(c')  \pmod{2} &\text{if } p = 2 \text{ and } f \ge 3. \end{cases}$ 
\end{prop}


\section*{Retained context: f=1 (source0.tex lines 277--283)}

\begin{Rem} \label{f=1}
For completeness, we recall that if $\chi$ and $\chi'$ are distinct characters on $\F_q^\times$ whose orders divide $p-1$, then $G(\chi) \ne G(\chi')$.  Clearly, $\chi$ and $\chi'$ are not primitive as soon as $f \ge 2$.   By \eqref{ConjChi}, we can assume that neither $\chi$ not $\chi'$ is trivial.  Define $n$ by $q-1 = (p-1)n$.  Since $\omega$ has order $q-1$ on $\F_q^\times$, we can write $\chi = \omega^{an}$ and $\chi' = \omega^{bn}$, where $1 \le a,b \le p-2$.  We have the $p$-adic digit expansion
$
an = a(1+p+\dots+ p^{f-1}), 
$
so $s_q(an) = af$ and similarly $s(bn) = bf$.  If $a$ and $b$ are not equal, Proposition \ref{EqualG} shows that $G(\chi) \ne G(\chi')$.
\end{Rem}


\section*{Retained context: gcd1 (source0.tex lines 355--357)}

\begin{prop} \label{gcd1}
Let $\chi$ and $\chi'$ be characters on $\F_q^\times$, with $\chi$ of order $q-1$.  If $G(\chi) = G(\chi')$, then $\chi$ and $\chi'$ are Frobenius-conjugate.
\end{prop}


\section{The counted unit: Proposition (source label Twist) and its complete original proof}

\begin{prop} \label{Twist}
Assume that $n  = \frac{q-1}{p-1}$ is prime and let $\chi$ and $\chi'$ be characters on $\F_q^\times$.  If $G(\chi \, \widehat{\eta}) = G(\chi' \,  \widehat{\eta})$ for all characters $\eta$ on $\F_p^\times$, then $\chi$ and $\chi'$ are Frobenius-conjugate.
\end{prop}

\begin{proof}
Suppose that $G(\chi \, \widehat{\eta}) = G(\chi' \,  \widehat{\eta})$ for all $\eta$.  If the order of both $\chi$ and $\chi'$ divides $p-1$, then $\chi = \chi'$ by Remark \ref{f=1}.  We may therefore assume that at least one of of these characters, say $\chi$, has the form $\chi = \omega^a$ with $n \nmid a$.  Also, $f \ge 2$, so $n \ge p+1$ and $n \nmid p-1$.    Take $\widehat{\eta} = \omega^{tn}$ with 
$$ 
t = \prod \{ \ell \text{ prime}, \, \ell \text{ divides } p-1 \text{ and } \ell \text{ does not divide } a \},
$$ 
allowing $t=1$ if the product is empty.  Let $\Psi = \chi \, \widehat{\eta} = \omega^c$ with $c = a + tn$ and observe that $\gcd(c,q-1) = 1$. Indeed, if $\ell$ is a prime dividing $q-1$, there are three cases. \vspace{2 pt}

\noindent \hspace{9 pt} i)  If $\ell = n$, then $\ell \nmid a$, so $\ell \nmid c$.  \hspace{7 pt} ii) If $\ell \vert p-1$ but $\ell  \nmid a$, then $\ell \vert t$, so $\ell \nmid c$.   \vspace{2 pt}

\noindent \hspace{3 pt} iii) If $\ell \vert p-1$ and $\ell \vert a$, then $\ell \nmid tn$, so $\ell \nmid c$. \vspace{2 pt}

\noindent Hence $\Psi$ has order $q-1$ and so, by Proposition \ref{gcd1}, $\chi' \, \widehat{\eta}$ is Frobenius-conjugate to $\chi \, \widehat{\eta}$.  But $\widehat{\eta}$ takes values in $\Mu_{p-1}$ so Frobenius at $p$ acts trivially on $\widehat{\eta}$.  It follows that $\chi$ and $\chi'$ are Frobenius-conjugate.
\end{proof}

\end{document}
